\begin{abstract}
We present a first-order behavioral simulation framework for organic optoelectronic compute-in-memory inference that links device-level assumptions to task-level behavior under quantization, variability, retention, analog-to-digital conversion, nonlinear write dynamics, and physical frontend effects. The framework is structured so that literature-anchored priors can later be replaced by measured device profiles without changing the training or evaluation code path. Across ResNet-18, ConvNeXt-Tiny, and Tiny-ViT-5M, the results show that nominal conductance quantization is not the dominant bottleneck once scale recovery and hardware-aware training are enabled. Instead, the practical limits depend strongly on task complexity, converter precision, hardware-instance transferability, proportional state-dependent noise, and nonlinear write dynamics. Cross-dataset experiments further show that the value of hardware-aware training grows sharply with task complexity but can become unstable in low-data regimes. We therefore position the framework not as a chip-predictive emulator, but as a calibrated materials-to-system bridge for determining which measured device characteristics most strongly constrain edge-vision deployment.
\end{abstract}
